\documentclass[10.5pt]{article}
\usepackage[letterpaper, margin=0.6in, top=0.5in, bottom=0.5in]{geometry}
\usepackage{amsmath,amssymb,amsthm}
\usepackage{fancyhdr}
\usepackage{xcolor}
\usepackage{tikz}
\usepackage{enumitem}
\usepackage{tcolorbox}
\usepackage{hyperref}

% Define colors matching instructor board annotations
\definecolor{chalkblue}{RGB}{0, 105, 185}
\definecolor{chalkred}{RGB}{205, 35, 35}
\definecolor{chalkpink}{RGB}{230, 90, 120}
\definecolor{chalkcyan}{RGB}{15, 160, 190}
\definecolor{chalkgreen}{RGB}{20, 130, 50}
\definecolor{boxbg}{RGB}{252, 252, 254}
\definecolor{boxframe}{RGB}{150, 160, 180}

\pagestyle{fancy}
\fancyhf{}
\renewcommand{\headrulewidth}{0pt}
\renewcommand{\footrulewidth}{0pt}

\setlength{\parindent}{0pt}
\setlength{\parskip}{2.5pt}

\begin{document}

% ==============================================================================
% PAGE 1
% ==============================================================================
\noindent \textbf{\large Math 105L} \hfill \textbf{\large Lecture 7-1}\\[1pt]
{\textbf{\Large Compound Interest \& Developing Exponential Growth}}\\[-5pt]
\rule{\textwidth}{1.2pt}

\vspace{0.03in}

\textbf{\large Preliminaries: Investing at $6\%$ Annual Interest}

Suppose you invest an initial principal $P_0 = \$10$ at $r = 6\%$ ($0.06$) annual interest, compounded annually:
\begin{itemize}[leftmargin=1.2em, topsep=1pt, itemsep=2pt]
  \item \textbf{After 1 year ($t=1$):}
  \[
    \$10 + 0.06 \times \$10 = \$10(1 + 0.06) = \mathbf{\$10.60}
  \]
  \item \textbf{After 2 years ($t=2$):}
  \[
    10.60 + 0.06 \times 10.60 = 10.60(1 + 0.06) = 10(1 + 0.06)(1 + 0.06) = \mathbf{10(1.06)^2 = \$11.236}
  \]
  \item \textbf{After 10 years ($t=10$):} $\mathbf{10(1.06)^{10} \approx \$17.91}$
  \item \textbf{After $t$ years:} $\mathbf{P(t) = 10(1.06)^t}$ \quad (Standard annual compounding: $P(t) = P_0 (1 + r)^t$)
\end{itemize}

\vspace{0.03in}

\textbf{\large Varying the Compounding Frequency ($n$ Periods Per Year)}

Suppose instead you invest $\$10$ at $6\%$ annual interest, compounded semi-annually ($n = 2$ times per year). That means each half-year you receive $\frac{6\%}{2} = 3\%$ interest:
\begin{itemize}[leftmargin=1.2em, topsep=1pt, itemsep=2pt]
  \item \textbf{After 1 year (2 periods):} $10\left(1 + \frac{0.06}{2}\right)^2 = 10(1.03)^2 = \$10.609 > \$10.60$.
  \item \textbf{After 2 years (4 periods):} $10\left(1 + \frac{0.06}{2}\right)^4 = 10(1.03)^4 \approx \$11.255$.
  \item \textbf{After 10 years (20 periods):} $10\left(1 + \frac{0.06}{2}\right)^{20} \approx \$18.06$.
  \item \textbf{General Compounding Formula ($n$ periods per year):}
  \[
    \mathbf{P(t) = P_0 \left(1 + \frac{r}{n}\right)^{nt}}
  \]
  \item \textbf{Why does more frequent compounding yield more money?}
  Because interest earned in earlier periods is added to the balance and itself earns interest in subsequent periods (\textit{interest on interest}).
\end{itemize}

\vspace{0.02in}

\textbf{\large Doubling the Principal vs.\ Doubling the Interest Rate}

\begin{itemize}[leftmargin=1.2em, topsep=1pt, itemsep=2pt]
  \item \textbf{Option A (Double Principal):} $\$20$ at $6\%$ annual interest $\implies P_A(t) = 20(1.06)^t$.
  \item \textbf{Option B (Double Rate):} $\$10$ at $12\%$ annual interest $\implies P_B(t) = 10(1.12)^t$.
  \item \textbf{Short-Term Horizon ($t = 10$ years):}
  \[
    P_A(10) = 20(1.06)^{10} \approx \mathbf{\$35.82}, \qquad P_B(10) = 10(1.12)^{10} \approx \mathbf{\$31.06} \quad \text{(Option A wins!)}
  \]
  \item \textbf{Long-Term Horizon ($t = 100$ years):}
  \[
    P_A(100) = 20(1.06)^{100} \approx \mathbf{\$6,786}, \qquad P_B(100) = 10(1.12)^{100} \approx \mathbf{\$835,220} \quad \text{(Option B crushes Option A!)}
  \]
  \item \textbf{Core Pedagogical Insight:} A linear multiplier on the initial principal ($P_0$) only scales linearly, whereas an increase in the exponential base compounds multiplicatively across time.
\end{itemize}

\newpage

% ==============================================================================
% PAGE 2
% ==============================================================================
\noindent \textbf{\large Math 105L} \hfill \textbf{\large Lecture 7-1 (Page 2)}\\[1pt]
{\textbf{\Large Pushing the Limit: Continuous Compounding \& The Number $e$}}\\[-5pt]
\rule{\textwidth}{1.2pt}

\vspace{0.03in}

\textbf{\large Pushing the Compounding Frequency to Infinity}

Suppose we invest $P_0 = \$1$ at $r = 100\%$ ($r = 1$) interest for $t = 1$ year. How much do we have if compounding occurs:
\begin{itemize}[leftmargin=1.2em, topsep=1pt, itemsep=2pt]
  \item \textbf{Annually ($n=1$):} $1\left(1 + \frac{1}{1}\right)^1 = \$2.00$
  \item \textbf{Semiannually ($n=2$):} $1\left(1 + \frac{1}{2}\right)^2 = (1.5)^2 = \$2.25$
  \item \textbf{Monthly ($n=12$):} $1\left(1 + \frac{1}{12}\right)^{12} \approx \mathbf{\$2.6130}$
  \item \textbf{Daily ($n=365$):} $1\left(1 + \frac{1}{365}\right)^{365} \approx \mathbf{\$2.7145}$
  \item \textbf{Every Minute ($n=525,600$):} $1\left(1 + \frac{1}{525,600}\right)^{525,600} \approx \mathbf{\$2.718279}$
  \item \textbf{Every Second ($n=31,536,000$):} $1\left(1 + \frac{1}{31,536,000}\right)^{31,536,000} \approx \mathbf{\$2.71828178}$
\end{itemize}

Does the balance grow without bound? \textbf{No!} It converges to a precise finite mathematical limit:

\begin{tcolorbox}[colback=chalkblue!6, colframe=chalkblue, boxrule=1pt, arc=3pt, left=8pt, right=8pt, top=4pt, bottom=4pt]
\textbf{\large Definition: Euler's Number $e$}
\vspace{1pt}

The mathematical constant $e$ is defined as the limit of continuous growth on unit capital at $100\%$ interest over one unit of time:
\[
  \mathbf{e = \lim_{n \to \infty} \left(1 + \frac{1}{n}\right)^n \approx 2.718281828459\dots}
\]
\end{tcolorbox}

\vspace{0.03in}

\textbf{\large General Continuous Compounding Formula}

For an arbitrary annual interest rate $r$ and principal $P_0$, we let $n \to \infty$ in the discrete compounding formula:
\begin{align*}
  \lim_{n \to \infty} P_0 \left(1 + \frac{r}{n}\right)^{nt} &= P_0 \lim_{n \to \infty} \left[\left(1 + \frac{1}{n/r}\right)^{n/r}\right]^{rt}
\end{align*}
Letting $m = \frac{n}{r}$, as $n \to \infty$, we have $m \to \infty$:
\[
  P_0 \left[\lim_{m \to \infty} \left(1 + \frac{1}{m}\right)^m\right]^{rt} = \mathbf{P_0 e^{rt}}
\]

\begin{tcolorbox}[colback=chalkgreen!6, colframe=chalkgreen, boxrule=1pt, arc=3pt, left=8pt, right=8pt, top=4pt, bottom=4pt]
\textbf{The Continuous Growth Law:}\\[2pt]
When compounding occurs \textit{continuously} ($n \to \infty$), the accumulated balance after $t$ years is:
\[
  \mathbf{P(t) = P_0 e^{rt}}
\]
\textbf{Example from Class:} Investing $\$1$ at $5\%$ ($r = 0.05$) compounded continuously for $1$ year yields:
\[
  P(1) = 1 \cdot e^{0.05 \cdot 1} = e^{0.05} \approx \mathbf{\$1.05127}
\]
\end{tcolorbox}

\end{document}

